A Hawkes process is a counting process whose intensity is driven by its own past.
Each event raises the intensity of future events by an amount that depends on the two types
involved and decays with the time elapsed.

\begin{defi}\label{def.hawkesprocess}
 A $\numEventTypes$-dimensional counting process $\multiCountingProc$ is called a Hawkes process
 if it admits an absolutely continuous compensator $\compensator$ with intensities
 \begin{equation}\label{eq.intensity_ordHawkes}
  \intensity(t) = \baseIntensity_e
  + \sum_{\eone=1}^{\numEventTypes}\intzerot \hawkesKernel\subscriptee(t-s) \, d\countingProc[\eone](s),
  \qquad e = 1,\dots,\numEventTypes,
 \end{equation}
 for some non-negative baselines $\baseIntensity_e \geq 0$,
 and some non-negative locally integrable functions $\hawkesKernel\subscriptee \geq 0$
 supported on the non-negative half line.
\end{defi}

The matrix-valued function
$t\mapsto [\hawkesKernel\subscriptee(t)]_{e,\eone = 1, \dots, \numEventTypes}$
is referred to as the kernel of the Hawkes process $\multiCountingProc$.
If all the kernel functions are integrable, we write
\begin{equation}\label{eq.branchingMatrix}
 \branchingMatrix\subscriptee := \Vert \hawkesKernel\subscriptee \Vert_{\Lone}
 = \int_{0}^{\infty} \hawkesKernel\subscriptee(t) \, dt
\end{equation}
for the matrix of $\Lone$ norms,
and we call the spectral radius $\branchingRatio$ of $\branchingMatrix$ the branching ratio of
the process;
if some of the kernel functions are not integrable, the branching ratio is set to $+\infty$.
Every statement of this chapter about stability is a statement about $\branchingRatio$,
and Section \ref{sec.stability} is devoted to reading it.
For the moment we record only that a $\numEventTypes$-dimensional Hawkes process admits a
stationary version with finite mean intensity if and only if $\branchingRatio < 1$.

\begin{remark}[the order of the indices]\label{remark.indexOrder}
 In \eqref{eq.intensity_ordHawkes} the influence that an event of type $\eone$ exerts on the
 intensity of type $e$ is written in the \emph{second} index:
 the row is excited and the column excites.
 The choice is not free of consequence.
 With this order, the sum in \eqref{eq.intensity_ordHawkes} is a matrix acting on a column
 vector, and so is every object we shall build from it;
 with the opposite order, each of them acquires a transpose.
 It also decides which marginal of $\branchingMatrix$ is the readable one.
 The \emph{column} sum $\sum_{e}\branchingMatrix\subscriptee$
 is the expected number of direct offspring of a single event of type $\eone$,
 a dimensionless count that answers ``how much does one event of this type beget'';
 the row sum answers no question anyone asks.

 The opposite order is common in the literature,
 and the two are hard to tell apart after the fact,
 because $\branchingRatio(\branchingMatrix) = \branchingRatio(\branchingMatrix^{\transpose})$:
 a transposed kernel passes every stability check,
 produces a stationary path of the right total rate,
 and is wrong only in which type excites which.
 It is a bug that still runs, and the only defence is to fix the convention once and state it.
\end{remark}

\subsection{The cluster representation}
\label{sec.clusterRepresentation}

Definition \ref{def.hawkesprocess} describes the process through its intensity.
There is a second description, due to \cite{HO74clu},
which describes the same process through a branching construction,
and which is the one that explains why the order flow carries information at all.

\begin{prop}[\cite{HO74clu}]\label{prop.clusterRepresentation}
 Assume $\hawkesKernel \geq 0$ and $\branchingRatio < 1$.
 Construct a point process as follows.
 \emph{i.} Immigrants of type $e$ arrive as a homogeneous Poisson process of rate
 $\baseIntensity_e$, independently across $e$.
 \emph{ii.} Every event of type $\eone$, immigrant or not,
 independently generates its direct offspring of type $e$
 as an inhomogeneous Poisson process on $(t,\infty)$ of intensity $\hawkesKernel\subscriptee(s-t)$,
 so that the number of such offspring is Poisson with mean $\branchingMatrix\subscriptee$.
 Then the superposition of all events so constructed is a stationary Hawkes process
 with baselines $\baseIntensity$ and kernel $\hawkesKernel$,
 and conversely every stationary Hawkes process with finite intensity admits this representation.
\end{prop}

The descendants of a single immigrant form a multitype Galton--Watson process
whose mean offspring matrix is $\branchingMatrix$,
and the population dies out almost surely, with finite expected total size,
precisely when $\branchingRatio < 1$.
The events of the process therefore partition into \emph{clusters},
each consisting of one immigrant and all of its descendants.

Two remarks on what a cluster is and is not.
A cluster is a run of events with a common ancestor.
It is not a run of events of the same type,
and in a market model it is not a run of same-signed events:
an offspring inherits its parent's type only to the extent that
$\branchingMatrix$ has mass on its diagonal blocks,
and in the parametrisations we use most of the mass is off them.
The distinction is measurable, and Section \ref{sec.readingOrderFlow} measures it.

The second remark concerns non-negativity.
That $\hawkesKernel \geq 0$ is required in Definition \ref{def.hawkesprocess}
is easy to read as a convenience, and it is not.
A negative entry has no reading as an offspring intensity,
so the branching construction of Proposition \ref{prop.clusterRepresentation}
does not exist;
the process one may still define by truncating the intensity at zero is no longer linear;
and $\branchingMatrix$, the cluster sizes, the endogenous fraction and the stationary intensity
of Section \ref{sec.stability} all lose their meaning together.
Every result in this chapter assumes it.

The cluster representation is also the answer to a question the rest of these notes will keep
asking: why should a sum of recent order flow predict anything?
Because a burst of events is, with high probability, a cluster in progress,
and a cluster in progress has descendants still to come.
A window sum is a crude estimator of ``a cluster is under way'',
and its predictive content is exactly the expected size of what remains.
